\section*{Aktivitas Pembelajaran dan Praktikum}
\addcontentsline{toc}{section}{Aktivitas Pembelajaran dan Praktikum}

\begin{table}[H]
\centering
\caption{Rencana aktivitas Bab V}
\label{tab:ch05-aktivitas}
\begin{tabularx}{\textwidth}{@{}llX@{}}
\toprule
\textbf{Sesi} & \textbf{Bentuk} & \textbf{Kegiatan} \\
\midrule
1 & Ceramah & Peran lexer dan antarmuka token stream \\
2 & Latihan & Menokenisasi program secara manual \\
3 & Diskusi & Ekspresi reguler dan prioritas aturan \\
4 & Praktikum & Menyusun spesifikasi token dan kasus uji \\
\bottomrule
\end{tabularx}
\end{table}

\subsection*{P1-A: Spesifikasi Token dan 30 Kasus Uji}
\textbf{Tujuan.} Menyusun spesifikasi token Subset-C yang lengkap
sebelum implementasi lexer.

\textbf{Langkah kerja.}
\begin{enumerate}[label=\arabic*.,leftmargin=*]
  \item Lengkapi \file{spec/tokens.md} dengan tabel pattern untuk
        setiap token.
  \item Tuliskan daftar kata kunci dan aturan escape untuk char serta
        string.
  \item Buat 30 kasus uji di \file{tests/lexer/}: 10 kasus normal,
        10 kasus batas, dan 10 kasus kesalahan.
  \item Untuk setiap kasus, tuliskan berkas \code{.sc} dan berkas
        \code{.expected} dengan format \code{BARIS:KOLOM JENIS LEXEME}.
  \item Periksa bahwa urutan aturan memenuhi maximal munch.
\end{enumerate}

\textbf{Definisi selesai.}
\begin{checklist}
  \item Semua token memiliki pattern.
  \item 30 kasus uji tersedia dan terdokumentasi.
  \item Format keluaran tes konsisten.
  \item Tidak ada aturan yang saling bertentangan.
\end{checklist}

P1-A belum dinilai secara sumatif. Hasilnya menjadi dasar Praktikum 1
(8\%) pada Bab VII.
